\end{lemma}

\begin{proof}
Now $b=4C-g$, and
\[
 9J(b)+1-2g=
 \begin{cases}
  2(54C+5),&g=0,\\
  2(54C-5),&g=1.
 \end{cases}
\]
The parenthesized integer is odd, so the valuation is exactly one.  Put
$z=3C-g$.  The shared exponent-two/one boundary gives $y=R(z)$.
When $g=0$, the integer $z$ is odd and
Lemma~\ref{lem:alt-factor} gives
\[
  3z+1=2^\lambda J(y).
\]
The returned source is $w=18C+1=2(3z+1)-1$, and therefore
\[
  w=Q_{\lambda+1}^{1}(J(y)).
\]
When $g=1$, the integer $z$ is even and the same lemma gives
\[
  3z+2=2^\lambda J(y).
\]
Now $w=18C-2=2(3z+2)$, so
\[
  w=Q_{\lambda+1}^{0}(J(y)).
\]
Here $\lambda\ge1$ is the maximal alternating-suffix length of $z$ when
$y>0$.  Thus in both phases
$w=Q_{\lambda+1}^{1-g}(J(y))$ and $m=\lambda+1\ge2$.  If $R(z)=0$,
Lemma~\ref{lem:alt-factor} instead gives a constant-tail state over
$J(0)=1$, which is exactly the source-zero alternative.
\end{proof}

Lemma~\ref{lem:D2} proves attachment of the first returned cursor to the
actual \LOSS\ token $y$.  It does not prove that every later raw and signed
exit preserves that attachment until a descendant is paid.  That latter
statement is the second principal global obligation.

\section{The exact remaining global obligations}

The earlier versions attempted to assemble local arithmetic into a marked
transition system with source anchors, one-shot token seeds, proof-token
multisets, and finite control counters.  The order-theoretic idea is sound,
but the semantic refinement to every legal \DRAW\ continuation is not yet
complete.  The missing content is now isolated as follows.

\begin{obligation}[Z: source-zero $j=1$ bootstrap]\label{obl:Z}
In the second branch of Lemma~\ref{lem:zero-j1}, where both the raw child and
its ordinary parent $y$ are \DRAW, construct a marked transition that either
pays a previously retained source/token component or consumes a declared
one-shot bootstrap which cannot be consumed again in the same earlier
fibre.  Prove that every later return to source zero is compatible with this
non-reseeding rule.
\end{obligation}

\begin{obligation}[H: high-return pair provenance]\label{obl:H}
Reserve an inner one-shot seed for the first long high-return pair
$\{u,c\}$, or prove that the pair consists of descendants of already retained
tokens.  After the seed has been consumed, prove for every subsequent long
high return in the same outer fibre that its incoming pair is exactly the
currently retained pair or certified descendants of it.  No merely finite
routing witness may be promoted to the rank.
\end{obligation}

\begin{obligation}[L: complete LOSS-anchored $D=2$ module]\label{obl:L}
Starting from the exact lift of Lemma~\ref{lem:D2}, enumerate every legal
raw/signed outcome branch, every remaining factor exponent, and the terminal
tail classes $m=1,2,3,\ge4$.  Prove that each macro-step either decreases an
explicit finite counter while retaining the same \LOSS\ token $y$, or
replaces $y$ by an actual child or certified descendant before any generic
obligation is re-entered.  In particular, prove that no branch exits as a
fresh free obligation at an unrelated arithmetic cursor.
\end{obligation}

\begin{obligation}[E: semantic exhaustiveness of the combined router]\label{obl:E}
After Obligations~\ref{obl:Z}--\ref{obl:L} are supplied, construct a
marked macro-relation satisfying every clause of
Definition~\ref{def:complete-router} below.  In particular, give a disjoint
and complete case partition for every control transition.  For each
equal-rank cycle, exhibit a natural counter that strictly decreases; for
each control reset, exhibit a preceding strict source or token replacement.
This proof must refine the actual game outcomes, not only an acyclic graph
of declared control labels.
\end{obligation}

\begin{remark}
Lemma~\ref{lem:fixed-tail-progress} closes the former fixed-tail part of the
entry countdown, and Lemma~\ref{lem:exp-one-progress} closes the local
exponent-one misuse.  They reduce Obligation~\ref{obl:E}, but do not imply
Obligations~\ref{obl:Z}--\ref{obl:L}.
\end{remark}

\section{A complete conditional closure theorem}

This section records the exact abstract statement supplied by the marked-rank
architecture.  The theorem is conditional, but the condition is now a
mathematically explicit property rather than the name of a control graph.

Let a marked configuration retain
\[
  \Pi=(s,\eta,\Phi(M),\zeta,\Psi(N),a),
\]
ordered lexicographically.  Here $s,a\in\Natz$ are genuine retained source anchors, $M,N$ are
finite multisets of certified token heights, $\Psi(N)$ is defined by the
same natural-sum formula as $\Phi(M)$, and $\eta,\zeta\in\{1,0\}$ are
one-shot seed bits ordered by $1>0$.  With the usual orders on the natural
and ordinal components, this finite lexicographic product is a well-order.
Temporary sources, phases, valuations, and tail lengths are not components
of $\Pi$ unless an explicit strict comparison installs them as retained
data.

\begin{definition}[Complete marked-router hypothesis]
\label{def:complete-router}
A marked macro-relation $\Longrightarrow$ is called complete for this game
when all of the following hold.
\begin{enumerate}[label=(C\arabic*)]
\item Every conjugated \DRAW\ state produces an initial marked configuration
      carrying that same actual \DRAW\ as its continuing cursor.
\item From every marked configuration carrying an actual \DRAW, a finite
      outcome-compatible normalization produces a successor marked
      configuration carrying an actual \DRAW.  Reverse-parent transformations
      are permitted only as proof transformations; whenever the cursor follows
      play, it follows a legal \DRAW\ child.
\item Every macro-step weakly decreases $\Pi$.  A token-changing step is a
      one-shot seed or a certified proper-descendant replacement; a
      source-changing step uses an explicit strict comparison of retained
      sources.
\item For each fixed value of $\Pi$, the restriction of
      $\Longrightarrow$ to that fibre is well-founded.  In particular, every
      equal-$\Pi$ control cycle has an explicitly decreasing natural counter.
\end{enumerate}
\end{definition}

\begin{lemma}[Well-founded fibre principle]\label{lem:fibre}
Let a transition relation carry a value in a well-order.  Suppose every
transition weakly decreases that value and the restriction of the relation
to each fixed-value fibre is well-founded.  Then the full relation is
well-founded.
\end{lemma}

\begin{proof}
An infinite path would induce a nonincreasing sequence in a well-order.  If
it had infinitely many strict decreases, those decreases would form an
infinite descending subsequence, which is impossible.  Hence its tail would
lie in one fixed fibre, contradicting fibre well-foundedness.
\end{proof}

\begin{theorem}[Conditional no-draw closure]\label{thm:conditional}
If a complete marked macro-relation in the sense of
Definition~\ref{def:complete-router} is constructed from the game relation,
then the conjugated game has no \DRAW\ positions.  Consequently every
original positive odd start has finite remoteness.
\end{theorem}

\begin{proof}
Assume a conjugated \DRAW\ exists.  Clause (C1) gives an initial marked
configuration carrying it.  Clause (C2), applied recursively, gives an
infinite sequence of marked configurations, each carrying an actual
continuing \DRAW.  Clauses (C3)--(C4) and
Lemma~\ref{lem:fibre} say that the same macro-relation is well-founded, so
such an infinite sequence cannot exist.  Hence no conjugated \DRAW\ exists.
Proposition~\ref{prop:first-normal} transfers the conclusion to every
original odd start: after the first move both possible children are
conjugated finite-remoteness states, so the original state is itself
\WIN\ or \LOSS\ with finite height.
\end{proof}

\begin{remark}
Theorem~\ref{thm:conditional} is not a proof of
Conjecture~\ref{conj:main}.  Obligations~\ref{obl:Z}--\ref{obl:L} are
specific missing semantic modules, and Obligation~\ref{obl:E} is the
remaining task of constructing a relation satisfying
Definition~\ref{def:complete-router}.  None is assumed silently in an
unconditional theorem.
\end{remark}

\section{Computational verification boundary}

The companion program \texttt{verify\_v7\_0.py} checks, over configurable
finite ranges:
\begin{enumerate}[label=(\arabic*)]
\item the original-to-binary normal form and strict contraction of $B$;
\item the long-tail recurrence and shared boundary;
\item the alternating-suffix factorization, the bijection $J$, ordinary
      coefficient-source recovery, the ordinary side relation, the
      source-boundary selector, and the exceptional child-pair identity;
\item the two permanent negative regressions;
\item every arithmetic identity in Lemmas
      \ref{lem:fixed-tail-progress}, \ref{lem:exp-one-progress},
      \ref{lem:zero-table}, \ref{lem:zero-j1},
      \ref{lem:high-return}, \ref{lem:D1}, and \ref{lem:D2}.
\end{enumerate}
The program does not infer infinite outcome classes and does not check
Obligations~\ref{obl:Z}--\ref{obl:E}.  Its final message explicitly says
that the run is supporting evidence, not a no-draw proof.

\section{Hostile-audit matrix}

\renewcommand{\arraystretch}{1.18}
\begin{longtable}{>{\raggedright\arraybackslash}p{0.24\textwidth}
                  >{\raggedright\arraybackslash}p{0.17\textwidth}
                  >{\raggedright\arraybackslash}p{0.49\textwidth}}
\toprule
Item & Status in v7.0 & Audit conclusion \\
\midrule
\endfirsthead
\toprule
Item & Status in v7.0 & Audit conclusion \\
\midrule
\endhead
Binary conjugation and $B(q)<q$ & Proved & Exact algebra; finite regression included. \\
Long-tail identities & Proved & Uniform in $a,r,e$; no bounded-modulus assumption. \\
Coefficient versus source & Corrected & The explicit $1\to3$ counterexample is retained permanently. \\
Multi-$A$ numerical anchor & Corrected & The $20\to30\to45\to68$, $B(68)=25$ counterexample is retained permanently. \\
Fixed-tail reverse entry & Proved locally & Lemma~\ref{lem:fixed-tail-progress} adds the missing continuing-DRAW and descendant-token alternatives. \\
Exponent-one reverse entry & Proved locally & Lemma~\ref{lem:exp-one-progress}; no use of the $r\ge2$ reverse lemma at $r=1$. \\
Source-zero $j=1$ row & Partly proved & Arithmetic and outcome split proved; global non-reseeding bootstrap is Obligation~\ref{obl:Z}. \\
Long high-return heights & Proved conditionally on token legitimacy & The inequalities are valid; rank insertion is isolated as Obligation~\ref{obl:H}. \\
Marked tail $D=1$ & Proved & Actual child of the retained LOSS token. \\
Marked tail $D=2$ & Exact first lift proved & Full raw/signed provenance remains Obligation~\ref{obl:L}. \\
Complete router & Open for the target theorem & No unconditional exhaustiveness claim; Obligation~\ref{obl:E}. \\
No-draw conclusion & Conditional only & Theorem~\ref{thm:conditional} under
Definition~\ref{def:complete-router}; Conjecture~\ref{conj:main} remains
open here. \\
\bottomrule
\end{longtable}

\section{Conclusion}

Version 7.0 does not repeat the cycle in which one local repair is followed
by a new hidden global assumption.  It makes three changes of principle.
First, it proves complete local progress statements for the fixed-tail and
exponent-one reverse rows rather than merely naming a finite witness.
Second, it never identifies a routing witness with a retained token without
a seed or a descendant comparison.  Third, it records the prize statement
as a conjecture until the source-zero bootstrap, high-return pair continuity,
LOSS-anchored $D=2$ module, and combined semantic exhaustiveness have all
been proved.

The result is a manuscript whose stated theorems are auditable and whose
remaining barrier is explicit.  It is not yet a solution of the prize
problem.  Any future Version 8 claiming the no-draw theorem must contain
proofs of Obligations~\ref{obl:Z}--\ref{obl:E}, or replace the marked-router
architecture by a genuinely different complete argument.

\section*{Acknowledgment of the audit}
Questions from Ingo Althoefer and subsequent hostile review exposed the
coefficient/source error, the exponent-one misuse, and the later token
provenance failures.  Version 7.0 treats those findings as permanent
regressions rather than cosmetic edits.

\begingroup
\raggedright
\begin{thebibliography}{9}
\bibitem{Guy1986}
R. K. Guy,
``John Isbell's Game of Beanstalk and John Conway's Game of Beans-Don't-Talk,''
\emph{Mathematics Magazine} 59(5) (1986), 259--269.
\href{https://doi.org/10.1080/0025570X.1986.11977258}{doi:10.1080/0025570X.1986.11977258}.

\bibitem{Guy1996}
R. K. Guy,
``Unsolved Problems in Combinatorial Games,'' Problem 42,
in \emph{Games of No Chance}, MSRI Publications 29,
Cambridge University Press, 1996.

\bibitem{AHZ2024}
I. Althoefer, M. Hartisch, and T. Zipproth,
``Analysis of a Collatz Game and Other Variants of the $3n+1$ Problem,''
in \emph{Advances in Computer Games}, LNCS 14528 (2024), 123--132.
\href{https://doi.org/10.1007/978-3-031-54968-7_11}{doi:10.1007/978-3-031-54968-7\_11}.

\bibitem{Prize}
I. Althoefer,
``Collatz Problems: The Two-Player $3n{+}{-}1$ Game,''
\url{https://althofer.de/collatz-prizes.html}, accessed 11 August 2026.

\bibitem{Zenodo}
G. Pochuev,
``Optimal Play in the Two-Player $3n\pm1$ Game,''
Zenodo, 2026.
\href{https://doi.org/10.5281/zenodo.21844684}{doi:10.5281/zenodo.21844684}.
The archived record predates the hostile audits summarized here.

\bibitem{Repo}
G. Pochuev,
``Optimal $3n\pm1$ Game: Proof and Verification Repository,'' 2026.
\url{https://github.com/Grisha-Pochuev/3n-plus-minus-1-game}.
\end{thebibliography}
\endgroup

\end{document}
